Out of hundreds of thousands of transcripts, the LLM was able to pull out a select few containing collusive content. This content takes a variety of forms including expressing a need for reduced industry capacity, noting that it has taken the lead in cutting supply and saying competitors need to do their part, and putting forth a call for the industry to show pricing discipline. At the same, there is a high rate of false positives. The LLM often did not catch the nuanced difference between a firm describing its own supply or capacity reduction, showing pricing discipline, or some other form of constrained conduct and the firm calling on the industry to do so.

While more exploratory work is needed, the evidence presented here supports a valuable role for an LLM in a screening process. An LLM can substitute for human auditors in performing an initial sifting through of a large corpus of transcripts. However, given the high rate of false positives and the critical importance of a competition agency avoiding the allocation of resources to non-collusive cases, human oversight of the LLM’s output is critical.

Here is a possible protocol:
\begin{description}
    \item[Step 1:] LLM reads and rates a corpus of transcripts.
    \item[Step 2:] For those transcripts with the highest scores, human auditors read the LLM-generated excerpts and identifies those with collusive content.
    \item[Step 3:] For those identified as having collusive conduct by the human audit, human auditors read the entire transcript and determines if it is worthy of further examination.
    \item[Step 4:] For those deemed worthy of further examination, all transcripts are collected for the company along with competitors’ transcripts and read by human auditors.
\end{description}

Based on the assessment of that body of transcripts, it can be determined if an official investigation should be initiated. As the LLM is improved in its performance, it will play more of a role by reducing the effort required in step 2 and assisting with step 4.
